\documentclass[11pt,a4paper]{article}

% =============================================================================
% Packages
% =============================================================================
\usepackage[margin=2.5cm]{geometry}
\usepackage{tabularx}
\usepackage{longtable}
\usepackage{booktabs}
\usepackage{xcolor}
\usepackage{listings}
\usepackage{hyperref}
\usepackage{amsmath}
\usepackage{amssymb}
\usepackage{enumitem}
\usepackage{parskip}
\usepackage[skins,breakable]{tcolorbox}

\definecolor{codebg}{rgb}{0.95,0.95,0.95}
\definecolor{notebg}{rgb}{0.93,0.96,0.99}
\definecolor{notebar}{rgb}{0.20,0.45,0.75}

% -----------------------------------------------------------------------------
% Code / inline macros
% -----------------------------------------------------------------------------
\newcommand{\code}[1]{\texttt{#1}}

% -----------------------------------------------------------------------------
% MATLAB code listing style
% -----------------------------------------------------------------------------
\lstdefinestyle{matlabstyle}{
  language=Matlab,
  backgroundcolor=\color{codebg},
  basicstyle=\ttfamily\small,
  breaklines=true,
  showstringspaces=false,
  frame=single,
  rulecolor=\color{gray!40},
  xleftmargin=4pt, xrightmargin=4pt,
  aboveskip=6pt, belowskip=6pt
}
\lstset{style=matlabstyle}

% -----------------------------------------------------------------------------
% Option/field table environments (plain tabular + booktabs), same pattern as
% the G09/G16 Toolbox Manual and the G_Utility Manual.
% -----------------------------------------------------------------------------
\newenvironment{argtable}[6]{%
  \begin{center}
  \renewcommand{\arraystretch}{1.25}
  \setlength{\dimen2}{\dimexpr\textwidth-#4-#5-6\tabcolsep-2\arrayrulewidth\relax}
  \begin{tabular}{@{}p{#4}p{#5}p{\dimen2}@{}}
  \toprule
  \textbf{#1} & \textbf{#2} & \textbf{#3} \\
  \midrule
}{%
  \bottomrule
  \end{tabular}
  \end{center}
}

\newenvironment{argtabletwo}[4]{%
  \begin{center}
  \renewcommand{\arraystretch}{1.25}
  \setlength{\dimen2}{\dimexpr\textwidth-#3-4\tabcolsep-2\arrayrulewidth\relax}
  \begin{tabular}{@{}p{#3}p{\dimen2}@{}}
  \toprule
  \textbf{#1} & \textbf{#2} \\
  \midrule
}{%
  \bottomrule
  \end{tabular}
  \end{center}
}

% -----------------------------------------------------------------------------
% Note / remark box
% -----------------------------------------------------------------------------
\newtcolorbox{tbnote}{
  enhanced, breakable,
  colback=notebg, colframe=notebar,
  boxrule=0pt, leftrule=2.5pt,
  arc=0pt, outer arc=0pt,
  left=8pt, right=8pt, top=6pt, bottom=6pt,
  before skip=8pt, after skip=8pt
}

\hypersetup{
  colorlinks=true,
  linkcolor=notebar,
  urlcolor=notebar,
  citecolor=notebar
}

% =============================================================================
\title{\textbf{RamanPCACluster}\\[4pt]
       \large An interactive app for unsupervised and supervised\\
       classification of Raman spectra}
\author{Sebastiano Trusso \\ CNR -- Istituto per i Processi Chimico-Fisici (IPCF), Messina, Italy}
\date{Version 1.0 --- 2026}

% =============================================================================
\begin{document}

\maketitle
\tableofcontents
\newpage

% =============================================================================
\section{Introduction}
% =============================================================================
\code{RamanPCACluster} is a self-contained MATLAB app (\code{PCAClusterApp.m}, with
a plain-script equivalent \code{PCA\_kmeans\_analysis.m} for batch/reproducible
runs) that takes a folder of Raman spectra and asks two related questions:

\begin{itemize}
  \item \textbf{Unsupervised:} does the data contain natural groupings at all,
  and if so, how many, and along which spectral features?
  \item \textbf{Supervised:} if the spectra also carry an externally known
  label (here, derived from each file's name -- e.g.\ the sampling site),
  how well does that particular grouping hold up statistically? Is it
  actually distinguishable from the spectra, or not?
\end{itemize}

The first question is answered by Principal Component Analysis (PCA),
$k$-means clustering, and Ward-linkage hierarchical clustering. The second is
answered by Linear Discriminant Analysis (LDA) with cross-validation. Both
paths share the same preprocessing (baseline removal, smoothing,
normalization) and the same dimensionality-reduced input (PCA scores), so
their results are directly comparable: a low Adjusted Rand Index between
$k$-means clusters and the known labels, \emph{together with} a low LDA
cross-validated accuracy, is good evidence that the known labels do not
correspond to a real chemical/spectral distinction in the data -- as opposed
to $k$-means simply having picked the wrong $k$ or the wrong features.

\begin{tbnote}
This manual documents the app's interface and gives a brief, self-contained
account of the theory behind each analysis step. It assumes no prior
familiarity with machine learning; readers already comfortable with PCA,
clustering, and LDA can skip directly to \S\ref{sec:gui} for the interface,
or \S\ref{sec:outputs} for the output file formats.
\end{tbnote}

% =============================================================================
\section{Requirements}
% =============================================================================
\begin{itemize}
  \item MATLAB with the \textbf{Statistics and Machine Learning Toolbox}
  (\code{pca}, \code{kmeans}, \code{silhouette}, \code{linkage},
  \code{dendrogram}, \code{fitcdiscr}, \code{crossval}).
  \item The \textbf{Signal Processing Toolbox} (\code{sgolayfilt}).
  \item No external dependencies beyond MATLAB itself: \code{backcor.m},
  \code{airPLS.m}, \code{snip.m}, \code{apls.m}, and \code{readdpt.m} are
  included alongside the app for self-containment.
\end{itemize}

% =============================================================================
\section{Data format}
% =============================================================================
The app reads every \code{.dpt} file in a chosen folder: two comma-separated
columns (Raman shift in cm$^{-1}$, intensity), no header. All spectra in the
folder must share the same wavenumber grid -- spectra are stacked directly
into a matrix rather than interpolated, so a file recorded on a different
grid raises an error instead of silently corrupting the comparison.

A group label is derived from each file name for later validation/plotting
only (it never influences the unsupervised steps): everything before the
underscore that introduces the sample identifier, itself always starting
with a digit -- e.g.\ \code{CT\_P\_11b.0.dpt} $\to$ \code{CT\_P},
\code{Paradiso\_31-nonpellet.0.dpt} $\to$ \code{Paradiso}. This is tolerant of
irregularly-formatted sample identifiers, which real-world file names
frequently contain.

% =============================================================================
\section{Getting started}
% =============================================================================
\begin{enumerate}
  \item Launch the app: \code{PCAClusterApp} at the MATLAB prompt.
  \item Click \textbf{Select spectra folder...} and choose the folder
  containing the \code{.dpt} files.
  \item Optionally narrow the \textbf{spectral range}, adjust
  \textbf{preprocessing} (baseline method and parameters, smoothing,
  normalization), and set the \textbf{number of clusters} $k$ (or let the
  app choose it automatically via the silhouette criterion).
  \item Click \textbf{Run analysis}. Results appear across the tabs on the
  right as each step completes.
  \item Click \textbf{Save results...} to export every plot as a PNG, the
  per-spectrum cluster assignments as a CSV, and all numeric results as a
  \code{.mat} file, into a folder of your choice.
\end{enumerate}

% =============================================================================
\section{Graphical user interface}
\label{sec:gui}
% =============================================================================
The window is split into a left sidebar (data selection and analysis
settings) and a tabbed results area on the right, plus a status bar along the
bottom that reports the current step and a one-line summary once the run
completes.

% -----------------------------------------------------------------------------
\subsection{Sidebar: Data}
% -----------------------------------------------------------------------------
\begin{argtabletwo}{Control}{Description}{3.2cm}{X}
\code{Select spectra folder...} & Opens a folder picker; loads every
\code{.dpt} file inside and reports the spectrum count. Shows
``Importing files...'' while reading, since a large folder can take a
moment. \\
\end{argtabletwo}

% -----------------------------------------------------------------------------
\subsection{Sidebar: Spectral range}
% -----------------------------------------------------------------------------
\begin{argtabletwo}{Control}{Description}{3.2cm}{X}
\code{Min} / \code{Max} & Restricts the analysis to a sub-range of the
wavenumber axis (cm$^{-1}$), e.g.\ to exclude a noisy edge or focus on a
region of chemical interest. \\
\code{Full range} & Resets Min/Max to the full wavenumber span of the
loaded spectra. \\
\end{argtabletwo}

% -----------------------------------------------------------------------------
\subsection{Sidebar: Preprocessing}
% -----------------------------------------------------------------------------
\begin{argtabletwo}{Control}{Description}{3.2cm}{X}
\code{Subtract baseline} & Enables/disables fluorescence-background
removal (\S\ref{sec:baseline}). \\
\code{Method} & Baseline algorithm: \code{backcor}, \code{airPLS},
\code{SNIP}, or \code{APLS} (\S\ref{sec:baseline}), each with its own
parameter fields shown below the dropdown. \\
\code{Savitzky-Golay smoothing} & Enables/disables local polynomial
smoothing (\S\ref{sec:smoothing}). \\
\code{Normalize} & \code{Area} (unit total intensity), \code{Max} (unit
peak height), \code{SNV} (standard normal variate), or \code{None}
(\S\ref{sec:normalize}). \\
\end{argtabletwo}

% -----------------------------------------------------------------------------
\subsection{Sidebar: Clustering}
% -----------------------------------------------------------------------------
\begin{argtabletwo}{Control}{Description}{3.2cm}{X}
\code{Choose k automatically} & Selects $k$ by maximizing mean silhouette
over $k=2,\dots,8$ (\S\ref{sec:choosek}), rather than using the fixed value
below. \\
\code{Number of clusters k} & Fixed number of $k$-means clusters, used
when automatic selection is off. \\
\code{Show confidence ellipses} & Overlays 80/85/90\% bivariate-normal
confidence ellipses (\S\ref{sec:ellipses}) on both PC1-PC2 scatter plots. \\
\code{Run analysis} & Executes the full pipeline (\S\ref{sec:pipeline})
and populates every results tab. \\
\end{argtabletwo}

% -----------------------------------------------------------------------------
\subsection{Sidebar: Reference library}
\label{sec:sidebar-references}
% -----------------------------------------------------------------------------
\begin{argtabletwo}{Control}{Description}{3.2cm}{X}
\code{Overlay reference spectra} & When checked, projects the loaded
reference library into the current PCA space and overlays it on both
PC1-PC2 scatter plots (\S\ref{sec:references}) at the next
\code{Run analysis}. Off by default. \\
\code{Select reference folder...} & Picks a folder of reference spectra
(flat \code{.txt} files, two columns, own native wavenumber grid per
file). A \code{SLoPP} folder next to the app, if present, is loaded
automatically at startup as the default. \\
\end{argtabletwo}

% -----------------------------------------------------------------------------
\subsection{Sidebar: Export}
% -----------------------------------------------------------------------------
\begin{argtabletwo}{Control}{Description}{3.2cm}{X}
\code{Save results...} & Exports every plot (PNG), the per-spectrum
cluster table (CSV), and all numeric results (\code{.mat}) to a chosen
folder (\S\ref{sec:outputs}). \\
\end{argtabletwo}

% -----------------------------------------------------------------------------
\subsection{Results tabs}
% -----------------------------------------------------------------------------
\begin{itemize}
  \item \textbf{Preprocessing} -- one example spectrum before/after
  baseline removal, smoothing, and normalization.
  \item \textbf{PCA} -- scree plot (per-component and cumulative explained
  variance).
  \item \textbf{Choose k} -- elbow plot (within-cluster sum of squares) and
  silhouette analysis over $k=2,\dots,8$.
  \item \textbf{Clusters} -- PC1-PC2 scatter, colored by filename-derived
  group (left) and by $k$-means cluster (right); optional confidence
  ellipses; optional reference-library overlay (\S\ref{sec:references}),
  one marker shape per material class.
  \item \textbf{Dendrogram} -- Ward-linkage hierarchical clustering tree.
  \item \textbf{LDA} -- canonical discriminant scores (LD1 vs.\ LD2) and
  the cross-validated confusion matrix.
  \item \textbf{GMM} -- PC1-PC2 scatter colored by the Gaussian Mixture
  Model's hard (argmax-posterior) cluster assignment, with each fitted
  component's own confidence ellipses; and a BIC-vs-number-of-components
  plot for model selection, with the number of components actually used
  (same $k$ as $k$-means) marked by a dashed line.
  \item \textbf{Loadings} -- PC1/PC2 loadings and per-cluster mean spectra,
  both plotted as an offset (waterfall) stack for readability.
\end{itemize}

% =============================================================================
\section{Analysis pipeline}
\label{sec:pipeline}
% =============================================================================
\begin{enumerate}
  \item \textbf{Load} every \code{.dpt} spectrum in the chosen folder onto a
  common wavenumber grid.
  \item \textbf{Restrict} to the chosen spectral range.
  \item \textbf{Preprocess}: baseline removal (optional, choice of method) +
  smoothing (optional) + normalization.
  \item \textbf{PCA}: reduce to the number of principal components needed to
  explain $\geq 95\%$ of the variance (capped at 10), used as the working
  representation for every step that follows.
  \item \textbf{Reference library} (optional, \S\ref{sec:references}):
  interpolates each reference spectrum onto the analysis grid, preprocesses
  it identically to the main dataset, and projects it into this same PCA
  space -- for visual comparison on the \textbf{Clusters} tab only, never
  fed into PCA/$k$-means/GMM/LDA. A reference whose own measured range does
  not cover the analysis range is skipped and reported.
  \item \textbf{Choose $k$}: elbow and silhouette analysis over
  $k = 2,\dots,8$.
  \item \textbf{$k$-means} on the retained PCA scores, with the fixed or
  silhouette-suggested $k$.
  \item \textbf{Compare} the resulting clusters against the filename-derived
  groups (contingency table and Adjusted Rand Index) -- informational only,
  never used to guide the unsupervised clustering itself.
  \item \textbf{Dendrogram}: Ward-linkage hierarchical clustering on the same
  retained PCA scores, as an alternative view that does not require choosing
  $k$ upfront.
  \item \textbf{LDA} (supervised, requires $\geq 2$ filename-derived groups):
  canonical discriminant scores plus cross-validated classification accuracy
  and confusion matrix, answering how separable the known groups actually
  are by any method -- not just how they happen to fall out of $k$-means.
  \item \textbf{GMM} (unsupervised, soft-clustering alternative): a Gaussian
  Mixture Model fitted by Expectation-Maximization on the same retained PCA
  scores, at the same $k$ as $k$-means, for a direct comparison; also
  reports the Adjusted Rand Index between the two partitions, and scans
  $k=1,\dots,8$ components to report each fit's BIC as an independent
  model-selection diagnostic.
  \item \textbf{Confidence ellipses} (optional): 80/85/90\% ellipses per
  group/cluster on the PCA scatter plots.
\end{enumerate}

\newpage
% =============================================================================
\section{Theoretical background}
% =============================================================================

% -----------------------------------------------------------------------------
\subsection{Baseline removal}
\label{sec:baseline}
% -----------------------------------------------------------------------------
Raman spectra of real samples typically sit on top of a smoothly-varying
fluorescence background. All four methods below estimate that background
$z$ from the raw spectrum $y$ and work by iteratively fitting a smooth curve
that is biased to stay \emph{below} the peaks -- the ways they achieve that
bias, and the smoothness they assume, are what distinguish them.

\paragraph{backcor \cite{Mazet2005}.} Fits a polynomial $z$ of a chosen
order to $y$, but instead of ordinary least squares it minimizes a
non-quadratic, asymmetric cost function of the residual $r = y - z$ that
penalizes positive residuals (points above the fit, i.e.\ peaks) more gently
than negative ones, so the fit is pulled down toward the lower envelope of
the spectrum. Four cost functions are available (\code{sh}, \code{ah},
\code{stq}, \code{atq}), differing in how sharply that asymmetry is applied;
the fit is refined iteratively until convergence.

\paragraph{airPLS \cite{Zhang2010}.} Adaptive Iteratively Reweighted
Penalized Least Squares: at each iteration, $z$ minimizes
\[
  \sum_i w_i (y_i - z_i)^2 \;+\; \lambda \sum_i \big(\Delta^d z\big)_i^2 ,
\]
a weighted residual term plus a roughness penalty on the $d$-th order
discrete derivative of $z$ (larger $\lambda$ forces a smoother baseline).
Weights $w_i$ start uniform and are then reduced for points where $y_i>z_i$
(candidate peaks), so subsequent iterations fit the baseline increasingly
around the peaks rather than through them.

\paragraph{SNIP \cite{Ryan1988}.} Statistics-sensitive Non-linear Iterative
Peak-clipping: starting from $y$, repeatedly applies
\[
  y_i \leftarrow \min\!\big(y_i,\; \tfrac{1}{2}(y_{i-M}+y_{i+M})\big)
\]
with a clipping window $M$ that is increased across iterations, which
progressively clips any point sitting above the local average of its
neighbors at that scale -- peaks get clipped down, smooth background does
not. An optional log-log-square-root (LLS) transform is applied before
clipping and inverted afterward, which compresses the dynamic range and
improves clipping of low-intensity peaks sitting on a high background.

\paragraph{APLS \cite{Cadusch2013}.} A further asymmetric penalized
least-squares variant tuned for fluorescence backgrounds under Poisson
(shot-noise) statistics, with its own asymmetric weighting scheme distinct
from airPLS's.

Only \code{airPLS} accepts the whole spectrum matrix in one call; the other
three are applied one spectrum at a time.

% -----------------------------------------------------------------------------
\subsection{Smoothing}
\label{sec:smoothing}
% -----------------------------------------------------------------------------
Savitzky-Golay smoothing (\code{sgolayfilt}) fits a low-order polynomial by
least squares within a sliding window centered on each point, and replaces
that point with the polynomial's value there. Unlike a moving average, it
tends to preserve peak height and width fairly well while still suppressing
high-frequency noise, which matters here since peak position/shape feeds
directly into the PCA that follows.

% -----------------------------------------------------------------------------
\subsection{Normalization}
\label{sec:normalize}
% -----------------------------------------------------------------------------
Removes spectrum-to-spectrum differences in overall signal level (laser
power, acquisition time, sample coupling) that are not chemically
meaningful, so that PCA is not simply driven by intensity scale:
\begin{itemize}
  \item \textbf{Area}: divide by $\int y\, d\tilde\nu$ (unit total
  intensity).
  \item \textbf{Max}: divide by $\max(y)$ (unit peak height).
  \item \textbf{SNV} (Standard Normal Variate): $(y-\bar y)/\mathrm{std}(y)$,
  zero mean and unit variance per spectrum.
  \item \textbf{None}: no normalization.
\end{itemize}

% -----------------------------------------------------------------------------
\subsection{Principal Component Analysis (PCA)}
\label{sec:pca}
% -----------------------------------------------------------------------------
PCA re-expresses each (preprocessed) spectrum, a point in a
high-dimensional space with one axis per wavenumber, in a new orthogonal
basis chosen so that the first axis (PC1) captures as much of the
spectrum-to-spectrum variance as possible, the second (PC2) captures as much
of the \emph{remaining} variance as possible while staying orthogonal to
PC1, and so on. Formally, with $X$ the (mean-centered) $n_{\text{spectra}}
\times n_{\text{points}}$ data matrix, PCA is the eigendecomposition of the
covariance matrix $\tfrac{1}{n-1}X^\top X = V\Lambda V^\top$: the columns of
$V$ are the \emph{loadings} (which wavenumbers drive each component), the
projections $\text{score} = XV$ are each spectrum's coordinates in the new
basis, and the eigenvalues $\Lambda$ give the variance -- hence the
\emph{explained variance percentage} -- captured by each component.

In this app, PCA serves as dimensionality reduction: only the leading
components needed to explain $\geq 95\%$ of the total variance (capped at
10) are kept for every downstream step ($k$-means, the dendrogram, LDA),
both to reduce noise from the long tail of high-order components and
because several of those downstream methods need more samples than
features to behave well.

% -----------------------------------------------------------------------------
\subsection{$k$-means clustering}
\label{sec:kmeans}
% -----------------------------------------------------------------------------
Given the PCA scores and a chosen number of clusters $k$, $k$-means
partitions the spectra into $k$ groups so as to minimize the total
within-cluster sum of squared distances to each cluster's own centroid,
\[
  \arg\min_{\{C_1,\dots,C_k\}} \sum_{j=1}^k \sum_{\mathbf{x}\in C_j}
  \lVert \mathbf{x} - \boldsymbol{\mu}_j \rVert^2 ,
\]
solved by Lloyd's algorithm (alternating between assigning each point to its
nearest current centroid, and recomputing centroids as the mean of their
assigned points) until convergence. Since Lloyd's algorithm can settle into
a poor local minimum depending on its random starting centroids, the app
runs many replicates (independent random restarts) and keeps the one with
the lowest total within-cluster distance. $k$-means is entirely
unsupervised: the filename-derived group labels play no role in it.

% -----------------------------------------------------------------------------
\subsection{Choosing $k$: elbow and silhouette}
\label{sec:choosek}
% -----------------------------------------------------------------------------
Since $k$-means requires $k$ as an input, two standard diagnostics are
computed over a range of candidate values ($k=2,\dots,8$):

\paragraph{Elbow plot.} The within-cluster sum of squares (WCSS) always
decreases as $k$ increases (more clusters can only fit the data at least as
well), but typically with diminishing returns past the ``true'' number of
groups -- the elbow, the point where adding another cluster stops buying
much improvement, is a visual (not automatic) heuristic for $k$.

\paragraph{Silhouette analysis.} For each point $i$, let $a(i)$ be its mean
distance to the other points in its own cluster (cohesion) and $b(i)$ its
mean distance to the points of the \emph{nearest other} cluster (separation);
its silhouette is
\[
  s(i) = \frac{b(i)-a(i)}{\max\big(a(i),\,b(i)\big)} \in [-1,1] ,
\]
close to $+1$ for a point that fits its own cluster well and far from
others, close to $0$ on a cluster boundary, and negative if it would fit
better in a different cluster. The app's ``choose $k$ automatically'' option
picks the $k$ that maximizes the mean silhouette across all points.

% -----------------------------------------------------------------------------
\subsection{Hierarchical clustering and dendrograms}
\label{sec:dendrogram}
% -----------------------------------------------------------------------------
Agglomerative hierarchical clustering builds a tree (dendrogram) bottom-up:
starting with every spectrum as its own cluster, it repeatedly merges the
two clusters whose merger increases the total within-cluster variance the
least (Ward's linkage criterion), until a single cluster remains. The
resulting tree height at which any two spectra are joined reflects how
dissimilar they are; cutting the tree at a given height yields a clustering
without having to commit to $k$ in advance, since every possible cut is
visible simultaneously. The app tunes its display's color threshold to
roughly match the $k$ chosen for $k$-means, purely so the two tabs tell a
visually consistent story, not because $k$ has any special meaning to the
dendrogram itself.

% -----------------------------------------------------------------------------
\subsection{Confidence ellipses}
\label{sec:ellipses}
% -----------------------------------------------------------------------------
For a set of 2D points (a group's or cluster's scores on PC1-PC2), assuming
they are drawn from a bivariate normal distribution with mean $\boldsymbol\mu$
and covariance $C$, a confidence ellipse at level $p$ is
\[
  \Big\{ \mathbf{x} : (\mathbf{x}-\boldsymbol\mu)^\top C^{-1}
  (\mathbf{x}-\boldsymbol\mu) \le \chi^2_2(p) \Big\} ,
\]
where $\chi^2_2(p)$ is the $p$-quantile of the chi-squared distribution with
2 degrees of freedom. In practice the ellipse is drawn from the eigenvectors
and eigenvalues of $C$ (its principal axes and their lengths), scaled by
$\sqrt{\chi^2_2(p)}$. The app draws the 80\%, 85\%, and 90\% ellipses (dotted,
dashed, solid) as a quick visual read on how tight or overlapping groups/
clusters are, beyond the raw scatter of points -- it is descriptive, not a
formal hypothesis test.

% -----------------------------------------------------------------------------
\subsection{Adjusted Rand Index}
\label{sec:ari}
% -----------------------------------------------------------------------------
The Adjusted Rand Index (ARI) measures agreement between two partitions of
the same $n$ items (here: the $k$-means clusters and the filename-derived
groups) by counting, over all $\binom{n}{2}$ pairs of items, how often the
two partitions agree on whether a pair is in the same group or not --
corrected for the level of agreement expected purely by chance:
\[
  \text{ARI} = \frac{\text{Index} - \text{Expected Index}}
                     {\text{Max Index} - \text{Expected Index}} ,
\]
computed from the contingency table between the two partitions. ARI $=1$
means the two partitions agree perfectly (up to a relabeling of cluster/group
names); ARI $\approx 0$ means the agreement is no better than a random
partition of the same sizes; ARI can go slightly negative for systematically
worse-than-random agreement. Since $k$-means never sees the filename-derived
labels, a low ARI does not mean the clustering failed -- it means the
groups implied by file name and the groups implied by spectral shape do not
coincide.

% -----------------------------------------------------------------------------
\subsection{Linear Discriminant Analysis (LDA)}
\label{sec:lda}
% -----------------------------------------------------------------------------
Unlike everything above, LDA is \emph{supervised}: it requires a known class
label per spectrum (here, the filename-derived group) and asks a different
question -- not ``what structure exists in the data on its own,'' but
``how well can these particular, already-known classes be told apart?''

\paragraph{Canonical discriminant scores.} With $C$ classes, define the
within-class scatter $S_w = \sum_{c=1}^{C}\sum_{\mathbf{x}\in
\text{class }c}(\mathbf{x}-\boldsymbol\mu_c)(\mathbf{x}-\boldsymbol\mu_c)^\top$
and the between-class scatter $S_b = \sum_{c=1}^C n_c
(\boldsymbol\mu_c-\boldsymbol\mu)(\boldsymbol\mu_c-\boldsymbol\mu)^\top$,
where $\boldsymbol\mu_c$ and $n_c$ are class $c$'s mean and size and
$\boldsymbol\mu$ the overall mean. LDA finds projection directions
$\mathbf{v}$ that maximize the ratio of between-class to within-class
variance,
\[
  \arg\max_{\mathbf v} \frac{\mathbf{v}^\top S_b \mathbf{v}}
                             {\mathbf{v}^\top S_w \mathbf{v}} ,
\]
solved as the generalized eigenvalue problem $S_b\mathbf{v} =
\lambda S_w \mathbf{v}$. There are at most $C-1$ non-trivial solutions
(discriminant axes LD1, LD2, \dots), ranked by eigenvalue -- the fraction of
between-class variance each one explains is reported analogously to PCA's
explained variance. Because $S_w$ is singular whenever there are more
features than samples (true of raw spectra almost always), this app runs
LDA on the same reduced PCA scores used for $k$-means, not on raw spectra.

\paragraph{Cross-validated classification accuracy.} The discriminant
scores above are useful for visualization, but a good-looking 2D separation
between classes does not by itself prove the classes are really
distinguishable (an LDA projection is chosen \emph{specifically} to
maximize apparent separation, so it can overstate how separable held-out
data would actually be). The app therefore also fits a discriminant
classifier (\code{fitcdiscr}) and evaluates it by $k$-fold cross-validation:
the data is split into $k$ folds, and for each fold in turn the classifier
is trained on the other $k-1$ folds and used to predict the held-out fold,
which it never saw during training. The fraction of correct predictions
across all folds is the cross-validated accuracy, and the resulting
predictions across all folds are assembled into a confusion matrix (rows:
true class; columns: predicted class), showing not just overall accuracy
but which specific classes are most often confused with which. The number
of folds is capped at the smallest class's own sample count, so no fold
ends up missing an entire class entirely.

\begin{tbnote}
\textbf{Reading ARI and LDA accuracy together.} A low ARI only says
$k$-means did not happen to reconstruct the known groups -- it leaves open
whether a smarter unsupervised method might, or whether the groups are
supervised-separable at all. A \emph{correspondingly} low cross-validated
LDA accuracy (i.e.\ not much better than $1/C$, chance level for $C$
balanced classes) closes that gap: it means that even a method built and
cross-validated \emph{specifically} to separate the known classes cannot do
so reliably, which is much stronger evidence that the known grouping does
not correspond to a real, consistent spectral/chemical distinction in the
data -- as opposed to merely being missed by $k$-means's particular choice
of $k$ or distance metric.
\end{tbnote}

% -----------------------------------------------------------------------------
\subsection{Gaussian Mixture Models (GMM)}
\label{sec:gmm}
% -----------------------------------------------------------------------------
Like $k$-means, a Gaussian Mixture Model is unsupervised -- it never sees
the filename-derived labels -- but it relaxes two assumptions $k$-means
makes implicitly. First, $k$-means assigns each point to exactly one
cluster; a GMM instead models the data as a weighted sum of $k$ Gaussian
components and assigns each point a full \emph{posterior probability} of
belonging to each one (\emph{soft} clustering). Second, $k$-means'
distance-to-centroid objective implicitly assumes every cluster is round
and similarly sized (equivalent to a GMM with equal, spherical, isotropic
covariances); a GMM instead fits each component's own mean
\emph{and} covariance matrix freely, so it can represent elongated or
correlated clusters that $k$-means would otherwise split or merge
incorrectly.

\paragraph{The mixture model.} The density of a point $\mathbf{x}$ (here, a
spectrum's retained PCA scores) is modeled as
\[
  p(\mathbf{x}) = \sum_{c=1}^k \pi_c\, \mathcal{N}(\mathbf{x};
  \boldsymbol\mu_c, \Sigma_c) ,
  \qquad \sum_{c=1}^k \pi_c = 1,\ \ \pi_c \ge 0 ,
\]
where $\mathcal{N}(\cdot;\boldsymbol\mu_c,\Sigma_c)$ is the multivariate
normal density of component $c$ and $\pi_c$ its mixing proportion (the
overall fraction of the data it accounts for). The parameters
$\{\pi_c,\boldsymbol\mu_c,\Sigma_c\}_{c=1}^k$ are fitted by maximum
likelihood via the Expectation-Maximization (EM) algorithm, alternating
between two steps until convergence:
\begin{itemize}[nosep]
  \item \textbf{E-step}: given the current parameters, compute each
  point's posterior (responsibility) for each component,
  \[
    \gamma_{ic} = P(c \mid \mathbf{x}_i) =
    \frac{\pi_c\, \mathcal{N}(\mathbf{x}_i;\boldsymbol\mu_c,\Sigma_c)}
         {\sum_{c'=1}^k \pi_{c'}\, \mathcal{N}(\mathbf{x}_i;\boldsymbol\mu_{c'},\Sigma_{c'})} ;
  \]
  \item \textbf{M-step}: given the responsibilities, re-estimate each
  component's mean, covariance, and mixing proportion as the
  responsibility-weighted mean/covariance/fraction of every point.
\end{itemize}
Like $k$-means' own alternating scheme, EM can converge to a local optimum
depending on its starting point, so the app runs several replicates and
keeps the highest-likelihood fit; a small regularization term is added to
each $\Sigma_c$ to keep it invertible when a component ends up with few or
nearly collinear points assigned to it. The app fits the GMM at the
\emph{same} $k$ as $k$-means (for a direct, like-for-like comparison) and
reports the two methods' Adjusted Rand Index (Sec.~\ref{sec:ari}) as a
measure of how much they actually agree, rather than assuming they do.

\paragraph{Hard assignment and confidence.} For visualization and for
comparison against $k$-means' own hard labels, each point is also assigned
to its most probable component, $\hat c_i = \arg\max_c \gamma_{ic}$; the
GMM tab reports the mean of $\max_c \gamma_{ic}$ across all points as a
one-number summary of how ``soft'' the fit actually turned out to be on
this data -- close to 1 means the components are cleanly separated (points
are confidently assigned), close to $1/k$ means many points sit
ambiguously between components. The ellipses drawn on the GMM scatter plot
use the same 80/85/90\% construction as Sec.~\ref{sec:ellipses}, but built
directly from each fitted component's own $\boldsymbol\mu_c,\Sigma_c$
(marginalized to PC1-PC2) rather than from the empirical covariance of
whichever points ended up assigned to it -- they show the \emph{model},
not a post-hoc summary of its output.

\paragraph{Model selection: BIC.} The number of components $k$ is not
estimated by the EM fit itself (a model with more components can always
fit the training data at least as well, exactly as WCSS with $k$-means).
The Bayesian Information Criterion trades off fit quality against model
complexity,
\[
  \mathrm{BIC} = -2\,\ell(\hat\theta) + p \ln n ,
\]
where $\ell(\hat\theta)$ is the fitted model's log-likelihood, $p$ its
number of free parameters (means, covariances, and mixing proportions all
grow with $k$), and $n$ the number of spectra -- lower BIC is better. The
app scans $k=1,\dots,8$ components and plots BIC against $k$ alongside the
GMM scatter, purely as an independent diagnostic (mirroring $k$-means' own
elbow/silhouette scan in Sec.~\ref{sec:choosek}); the $k$ actually used for
the fitted model shown is still tied to $k$-means' choice, not
automatically reset to BIC's minimum.

% -----------------------------------------------------------------------------
\subsection{Reference-library overlay (out-of-sample PCA projection)}
\label{sec:references}
% -----------------------------------------------------------------------------
Every method above analyzes the loaded spectra as a whole: PCA is fitted on
them, and $k$-means/GMM/LDA all partition or classify them. A
\emph{reference library} -- e.g.\ SLoPP/SLoPP-E, a published collection of
Raman spectra of known plastic/polymer materials \cite{Munno2020} -- serves
a different purpose: comparing the actual samples against known standards,
without those standards influencing the analysis in any way. This is
\emph{out-of-sample projection}: mapping new points into an
\emph{already-fitted} PCA space, never refitting that space to include them.

\paragraph{Why not just add them to the dataset.} Including reference
spectra directly in the PCA/clustering input would let them pull the
principal axes and cluster boundaries toward themselves -- the axes would
then reflect a mix of ``how the actual samples vary'' and ``how the actual
samples differ from a large, possibly unrelated, standards library'',
which is not the question being asked. Keeping them out of the fit and
only projecting them \emph{afterward} guarantees the PCA space reflects
the samples alone, with the references mapped into it purely for visual
reading.

\paragraph{The projection.} PCA's own \textsc{score} output is defined as
$\text{SCORE} = (X - \bar{\mathbf{x}})\,\text{COEFF}$, where $\bar{\mathbf x}$
is the analyzed dataset's own mean (row-wise) and \textsc{coeff} its
loadings (principal axes). A reference spectrum $\mathbf{x}_{\text{ref}}$,
once preprocessed \emph{identically} to the main dataset (same baseline
method, smoothing, normalization -- otherwise ``the same PCA space'' would
not mean the same thing for the two data sets) and interpolated onto the
exact wavenumber grid the analysis was run on, is mapped into the same
space by reusing that same mean and those same loadings,
\[
  \text{score}_{\text{ref}} = (\mathbf{x}_{\text{ref}} - \bar{\mathbf{x}})\,\text{COEFF} ,
\]
\emph{without} re-running \code{pca} on it. Since each reference spectrum
in a library assembled from many instruments/sessions is generally
recorded on its own native wavenumber grid (unlike the main dataset, which
must already share one), it is linearly interpolated onto the analysis
grid first; a reference whose own measured range does not fully cover the
analysis range is left out rather than extrapolated, since a baseline or
peak beyond what was actually measured for that reference cannot be
inferred reliably.

\paragraph{Display.} Each reference spectrum's file name is parsed for its
material class (e.g.\ ``Polyethylene 12. Blue Fragment.txt'' $\to$
``Polyethylene''), and every class is drawn with its own marker shape
(cycling through a fixed set) and color, with a legend entry -- visually
distinct from the filename-derived groups' or $k$-means clusters' own
markers/colors, so a reference point is never mistaken for an actual
sample.

% =============================================================================
\section{Output files}
\label{sec:outputs}
% =============================================================================
\textbf{Save results...} exports, into the chosen folder:

\begin{itemize}
  \item \code{preprocessing\_raw\_and\_baseline.png},
  \code{preprocessing\_processed.png} -- example spectrum before/after
  preprocessing.
  \item \code{scree\_plot.png} -- PCA explained variance.
  \item \code{elbow\_plot.png}, \code{silhouette\_plot.png} -- $k$ selection
  diagnostics.
  \item \code{clusters\_by\_filename\_group.png},
  \code{clusters\_by\_kmeans.png} -- PC1-PC2 scatter plots.
  \item \code{dendrogram.png} -- hierarchical clustering tree.
  \item \code{lda\_scatter.png}, \code{lda\_confusion\_matrix.png} -- LDA
  results.
  \item \code{gmm\_scatter.png}, \code{gmm\_bic.png} -- GMM soft-clustering
  scatter (with component ellipses) and the BIC model-selection scan.
  \item \code{pca\_loadings.png}, \code{mean\_spectrum\_per\_cluster.png} --
  waterfall-stacked loadings and per-cluster mean spectra.
  \item \code{cluster\_assignments.csv} -- one row per spectrum: file name,
  filename-derived group, $k$-means cluster, GMM cluster (hard assignment),
  and PC1/PC2/PC3 scores.
  \item \code{reference\_projections.csv} (only if the reference overlay
  was used) -- one row per reference spectrum actually shown: name,
  material class, and its projected PC1/PC2/PC3 scores.
  \item \code{pca\_kmeans\_results.mat} -- every numeric result (raw and
  preprocessed data, PCA coefficients/scores/explained variance, cluster
  assignments, contingency table, Adjusted Rand Index, elbow/silhouette
  curves, the full LDA output: discriminant scores, explained variance,
  cross-validated accuracy, confusion matrix, class names; the full GMM
  output: the fitted \code{gmdistribution} model, hard cluster assignment,
  posterior probabilities, BIC, the BIC scan across $k=1,\dots,8$, and the
  Adjusted Rand Index against both $k$-means and the filename-derived
  groups; and, if used, the reference library's folder, projected scores,
  material classes, names, and the list of references skipped for not
  covering the analysis range).
\end{itemize}

% =============================================================================
\begin{thebibliography}{9}
% =============================================================================
\bibitem{Mazet2005}
V. Mazet, C. Carteret, D. Brie, J. Idier, B. Humbert, ``Background removal
from spectra by designing and minimising a non-quadratic cost function,''
\emph{Chemometrics and Intelligent Laboratory Systems}, 76(2), 121--133
(2005).

\bibitem{Zhang2010}
Z.-M. Zhang, S. Chen, Y.-Z. Liang, ``Baseline correction using adaptive
iteratively reweighted penalized least squares,'' \emph{Analyst}, 135(5),
1138--1146 (2010).

\bibitem{Ryan1988}
C. G. Ryan, E. Clayton, W. L. Griffin, S. H. Sie, D. R. Cousens, ``SNIP, a
statistics-sensitive background treatment for the quantitative analysis of
PIXE spectra in geoscience applications,'' \emph{Nuclear Instruments and
Methods in Physics Research B}, 34, 396--402 (1988).

\bibitem{Cadusch2013}
P. J. Cadusch, M. M. Hlaing, S. A. Wade, S. L. McArthur, P. R. Stoddart,
``Improved methods for fluorescence background subtraction from Raman
spectra,'' \emph{Journal of Raman Spectroscopy}, 44(11), 1587--1595 (2013).

\bibitem{Munno2020}
K. Munno, H. De Frond, B. O'Donnell, C. M. Rochman, ``Increasing the
accessibility for characterizing microplastics: introducing new
application-based and spectral libraries of plastic particles (SLoPP and
SLoPP-E),'' \emph{Analytical Chemistry}, 92(3), 2443--2451 (2020).
\emph{[Citation as recalled -- please verify page/DOI before relying on
it in a publication.]}

\end{thebibliography}

\vspace{1cm}
\noindent\textit{Developed with the assistance of an AI coding tool (Claude,
Anthropic), under the author's supervision and review.}

\end{document}
